\documentclass[10pt,a4paper]{amsart}
\usepackage[margin=19mm]{geometry}
\usepackage{amsmath,amssymb,mathtools}
\usepackage[scr=rsfs,cal=euler]{mathalfa}
\usepackage{xfrac}
\usepackage[unicode,hidelinks,bookmarks=false]{hyperref}
\newcommand{\ie}{i.e.\ }
\newcommand{\cf}{cf.\ }
\newcommand{\resp}{respectively\ }
\newcommand{\Subsection}{\subsection}
\providecommand{\nobreakdash}{\nobreak\hbox{-}}
\setlength{\emergencystretch}{2em}
\multlinegap0pt
\usepackage{tikz}
\usepackage{amssymb}
\usepackage{mathtools}
\newtheorem{proposition}{Proposition}
\def\ZD{\ensuremath{\mathscr D}}
\def\ZN{\ensuremath{\mathbb N}}
\def\ZT{\ensuremath{\mathbb T}}
\newcommand {\e }[1]{\eqref{#1}}
\title{On UC-multipliers for multiple trigonometric systems}
\author{Grigori A. Karagulyan}
\date{Source: arXiv:2601.10360v1 (CC BY 4.0)}
\begin{document}
\maketitle

\noindent\emph{Source and licence.} On UC-multipliers for multiple trigonometric systems. Version arXiv:2601.10360v1 (CC BY 4.0), distributed by arXiv under the Creative Commons Attribution 4.0 International licence, http://creativecommons.org/licenses/by/4.0/, as recorded in the arXiv licence metadata for this exact version and shown on the abstract page https://arxiv.org/abs/2601.10360v1. Full licence notice: \texttt{licenses/arxiv-2601.10360-CC-BY-4.0.txt}.\par\smallskip

\noindent\textit{Editorial scope.} Counted unit: the article's proposition P1 (source0.tex lines 328--330). The article's complete original proof of it is delivered with it (source0.tex lines 331--366, one \texttt{proof} environment of 36 lines). No mathematics is written, repaired or re-derived: the statement and the proof are the article's own lines, in the article's own notation, and no retained line is substituted or re-flowed.  No further source-local context is needed: every cross-reference of every retained line resolves inside the two delivered spans.  Nothing was omitted from any retained span, so the package carries no omission record.  No retained line contains a citation, so the wrapper carries no bibliography and no bibliography entry.  No retained line contains a figure environment, an image-inclusion command or drawing code, so no image is delivered and the package declares no asset dependency.  Editorial formatting only, in addition: an AMS \texttt{amsart} preamble that restates the article's own declarations and macros used by the retained lines, namely the environments \texttt{proposition} on separate counters, together with the standard packages those lines need.  The retained environments print in this document as Proposition 1 in order of appearance, whereas the article prints the same items with the numbering of its own full document.  The complete native source of the article as distributed in the arXiv e-print for this exact version is bundled unchanged as \texttt{sources/192757/source0.tex}.
\begin{proposition}\label{P1}
	If $p_1,\ldots, p_d$ are mutually coprime numbers and $p=p_1\ldots p_d$, then the  systems $\ZD^{(p)}$ and $\ZD^{(p_1,\ldots,p_d)}$ are probabilistically equivalent. Moreover, this equivalence is generated by a measure-preserving map $\Theta:\ZT^p\to \ZT$.
\end{proposition}

\begin{proof}
According to the Chinese remainder theorem for any integers  $n_j\in \ZN_{p_j}$, $j=1,2,\ldots,d$ there is a unique integer $l\in \ZN_{p}$ such that
\begin{align*}
	l=n_j\mod p_j,\quad j=1,2,\ldots,d.
\end{align*}
This defines a one-to-one mapping  $\tau:\ZN_{p_1,\ldots,p_d}\to \ZN_{p}$ such that $\tau(n_1,n_2,\ldots, n_d)=l$. Consider the mapping $\bar \tau:\sqcap_{j=1}^d\ZN_{p_j}\to \ZN_{p}$ defined by
\begin{equation*}
	\bar\tau(u_1,\ldots,u_d)=\left\langle\tau(u_1,\ldots,u_d)\sum_{j=1}^d\frac{p}{p_j}\right\rangle_p.
\end{equation*}
Observe that this is a one-to-one mapping. 
Since the sets $\sqcap_{j=1}^d\ZN_{p_j}$ and $\ZN_{p}$ both have cardinality $p$ it remains to show that  $\bar \tau$ is injective. For $ (u_1,\ldots,u_d)\neq (u'_1,\ldots,u'_d)$ we have
	\begin{align}
		&\left\langle\bar \tau (u_1,\ldots,u_d)-\bar \tau(u'_1,\ldots,u'_d)\right\rangle_p\\
		&\qquad\qquad\qquad\qquad =\left\langle	(\tau (u_1,\ldots,u_d)- \tau(u'_1,\ldots,u'_d))\sum_{j=1}^d\frac{p}{p_j}\right\rangle_p.\label{x1}
	\end{align}
	One can check that $\tau (u_1,\ldots,u_d)- \tau(u'_1,\ldots,u'_d)$ is not divisible by $p$, as well as  $\sum_{j=1}^d\frac{p}{p_j}$ and $p$ are coprime numbers. This implies that \e{x1} can not be $0$ and so $\bar \tau$ is injective. Using $\bar\tau$ we define a MP-map $\Theta:\ZT^d\to \ZT$ by choosing 
	\begin{equation*}
		\Theta(\delta_{u_1}^{(p_1)}\times \cdots\times \delta_{u_d}^{(p_d)})=\delta^{(p)}_u \text{ if } u=\bar \tau(u_1,\ldots,u_d),
	\end{equation*}
while inside of each $\delta_{u_1}^{(p_1)}\times \cdots\times \delta_{u_d}^{(p_d)}$ the map $\Theta$ is defined arbitrarily,  provided that the measure-preserving property is maintained. Using the definitions of $\tau$ and $\bar \tau$, we obtain
\begin{align*}
	&\left\{\frac{\tau(n_1,\ldots,n_d)\bar \tau(u_1,\ldots,u_d)}{p}\right\}\\
	&\qquad\qquad\qquad\qquad=	\left\{\frac{\tau(n_1,\ldots,n_d)\tau(u_1,\ldots,u_d)\sum_{j=1}^d\frac{p}{p_j}}{p}\right\}\\
	&\qquad\qquad\qquad\qquad=	\left\{\sum_{j=1}^d\frac{\tau(n_1,\ldots,n_d)\tau(u_1,\ldots,u_d)}{p_j}\right\}=\left\{\sum_{j=1}^d\frac{n_ju_j}{p_j}\right\}.
\end{align*}
Thus for $n=\tau(n_1,\ldots,n_d)$ and $u=\bar \tau (u_1,\ldots,u_d)$ we have
\begin{align*}
	\exp\left(2\pi i\cdot  \frac{nu}{p}\right)&=\exp\left(2\pi i\cdot  \frac{\bar\tau(n_1,\ldots,n_d) \tau(u_1,\ldots,u_d)}{p}\right)\\
	&=\prod_{j=1}^d\exp\left(2\pi i\cdot  \frac{n_j u_j}{p_j}\right).
\end{align*}
This implies that 
\begin{equation*}
	\Theta(t_{n_1,\ldots,n_d}^{(p_1,\ldots,p_d)})=t^{(p)}_n,\text{ whenever } n=\tau(n_1,\ldots,n_d). 
\end{equation*}
Therefore the systems $\ZD^{(p)}$ and $\ZD^{(p_1,\ldots,p_d)}$ are probabilistically equivalent.
\end{proof}

\end{document}
